\section{Related Work}
\label{sec:related}

Prioritizing simulation-based tests for self-driving cars (SDCs) has been investigated from three main paradigms: (i)~\emph{search-based and diversity} methods that reorder test suites to maximize behavioral or static distance, (ii)~\emph{feature-based machine learning} that scores roads via hand-engineered tabular statistics, and (iii)~\emph{deep sequence learning} that processes raw geometry end-to-end. 
In this section, we review each paradigm, analyze the theoretical foundations of group equivariance and physical regularization, and position SE2RoadNet within the literature.

\subsection{Search-Based and Diversity-Driven Prioritization}

SDC-Prioritizer~\cite{sdcprio} formulates test case prioritization as a black-box combinatorial optimization problem.
A chromosome $\pi = \langle t_1, \dots, t_n \rangle$ represents a candidate execution order, evolved via genetic algorithms to ensure that diverse and computationally inexpensive tests are scheduled first.
The single-objective variant maximizes a cost-discounted diversity metric:
\begin{equation}
f(\pi) = \sum_{i=2}^{n} \frac{\mathrm{dist}(t_i, t_{i-1})}{\mathrm{cost}(t_i) \cdot i},
\end{equation}
where $\mathrm{dist}(t_i, t_{i-1})$ denotes the Euclidean distance between $z$-normalized static road feature vectors, and the multiplier $i$ prioritizes early fault detection.
The multi-objective variant (MO-SDC-Prioritizer) uses NSGA-II to optimize diversity and execution cost simultaneously.
While search-based prioritization requires no historical failure labels and operates with minimal runtime overhead ($<0.5\%$), diversity serves only as an indirect surrogate for fault-revealing capability.
Crucially, global Euclidean distance over hand-crafted features provides neither rotational invariance nor sampling-rate consistency.

\subsection{Feature-Based Machine Learning}

SDC-Scissor~\cite{sdcscissor,saner} predicts unsafe scenarios prior to simulation by extracting static descriptors---such as turn counts, min/max/mean turn angles, pivot radii, total length, and Shapely polygon bounding areas---and feeding them into tabular classifiers (Logistic Regression, Random Forest, Gradient Boosting).
Although SDC-Scissor reports cost-effectiveness gains of up to $4.0\times$ by filtering out benign scenarios, it suffers from two structural weaknesses:
First, counting features (e.g., number of left/right turns) depend on an arbitrary spline segmentation threshold; resampling the road alters these descriptors.
Second, aggregating roads into scalar summaries erases the localized spatial transitions where sudden steering adjustments induce vehicle spinouts.

\subsection{Deep Sequence Models}

To capture fine-grained geometry, recent research has turned to deep neural networks.
ITS4SDC~\cite{its4sdc} models roads as sequences of per-segment relative turn angles and lengths $(\Delta\theta_k, \ell_k)$ using a bi-directional LSTM (220 units), reaching an $F_1$-score of $0.89$ in the ICST 2025 competition.
Setting the initial heading to zero partially mitigates global orientation bias, but remains a heuristic rather than an architectural invariance guarantee.
Furthermore, the LSTM's rigid sequence dimension ($N=197$) ties predictions to a specific discretization rate.

Our earlier baseline, RoadFury~\cite{roadfury}, introduced a 4-layer Pre-LN Transformer Encoder~\cite{transformer,preln} with Stochastic Weight Averaging (SWA)~\cite{swa} trained on a 10-channel geometric matrix.
While RoadFury won the ICST 2026 competition with an APFD of $0.804 \pm 0.012$, its feature set retained frame-dependent components ($\sin\theta, \cos\theta, \theta$) and absolute sinusoidal positional embeddings, making it vulnerable to coordinate rotation and scale shifts.

\subsection{Theoretical Foundations: Invariance and Physics Grounding}

\textbf{Group Equivariance}: 
A mapping $\Phi: \mathcal{X} \to \mathcal{Y}$ is equivariant to a transformation group $G$ if $\Phi(T_g x) = T'_g \Phi(x)$ for all $g \in G$~\cite{gcnn}.
When the target space is invariant ($T'_g = \mathrm{Id}$), the output satisfies $\Phi(T_g x) = \Phi(x)$.
Standard Group Equivariant CNNs~\cite{gcnn} operate over discrete planar rotations ($p4$ or $p4m$ symmetries on 2D pixel grids) and cannot provide continuous $\mathrm{SE}(2)$ rigid-motion invariance for continuous 1D spatial curves.
In contrast, by the fundamental theorem of planar curves~\cite{frenet}, a regular curve is uniquely determined up to a rigid motion by its signed curvature $\kappa(s)$ parameterized by arclength $s$.
By designing an attention mechanism conditioned solely on intrinsic differential curvature and relative arclength intervals $\Delta s = |s_i - s_j|$ via Random Fourier Features (RFF)~\cite{rff}, SE2RoadNet achieves continuous, exact $\mathrm{SE}(2)$ invariance.

\textbf{Physics-Informed Regularization}:
In autonomous vehicle dynamics, lateral wheel slip occurs when centripetal acceleration exceeds available tire-road friction: $a_c = v^2 \kappa(s) > \mu g$.
Physics-Informed Neural Networks (PINN)~\cite{pinn} integrate physical conservation laws into neural loss functions.
In this work, we embed the physical friction limit directly into the prioritization objective, penalizing predictions where high-risk scores are assigned to low-curvature trajectories.

\subsection{Positioning}

Table~\ref{tab:positioning} contrasts SE2RoadNet against representative prior approaches.
While empirical APFD has largely plateaued around $0.80$, SE2RoadNet is the first framework to provide both architectural $\mathrm{SE}(2)$ invariance and physics-grounded curvature monotonicity.

\begin{table}[t]
\caption{Comparison of SDC Test Prioritization Approaches.}
\label{tab:positioning}
\centering{\small
\setlength{\tabcolsep}{4pt}
\begin{tabular}{llccc}
\toprule
\textbf{Method} & \textbf{Paradigm} & \textbf{APFD} & \textbf{$\mathrm{SE}(2)$-Inv.} & \textbf{Physics Loss} \\
\midrule
Random & Baseline & $\approx 0.50$ & \xmark & \xmark \\
SO-SDC-Prioritizer~\cite{sdcprio} & Search & $\approx 0.77$ & \xmark & \xmark \\
Greedy Diversity~\cite{sdcprio} & Search & $\approx 0.80$ & \xmark & \xmark \\
SDC-Scissor~\cite{sdcscissor} & Tabular ML & $\approx 0.75$ & \xmark & \xmark \\
ITS4SDC~\cite{its4sdc} & Deep (LSTM) & $\approx 0.78$ & \xmark & \xmark \\
RoadFury~\cite{roadfury} & Deep (Transf.) & $\approx 0.81$ & \xmark & \xmark \\
\textbf{SE2RoadNet (Ours)} & \textbf{Deep (Transf.)} & \textbf{$\approx 0.81$} & \textbf{\cmark} & \textbf{\cmark} \\
\bottomrule
\end{tabular}}
\end{table}
